\section*{Part II: A Decision Architecture for All Specialized Roles}

\subsection*{9. Design a Scorecard Around Outcomes, Not Adjectival Virtues}

Build each role from \textbf{3 to 6 competencies}, not 25 traits. For each competency specify observable behavior, an assessment source, and a minimum acceptable level. Separate \textbf{gates} (must be met) from \textbf{tradeable dimensions} (a strength may offset a weakness). For a reviewer, correctness and honesty about uncertainty might be gates, while presentation polish can be coached. For a sales leader, ethical representation and accurate pipeline judgment are gates; familiarity with the firm's CRM is trainable.

Every criterion should answer, ``What would a strong person \textit{do} differently in a real situation?'' Write \textit{detects data leakage before praising the model}, not \textit{has strong AI judgment}. Write \textit{names the budget owner and explains the approval path from a real deal}, not \textit{commercially savvy}. The more specific the behavioral anchor, the harder it becomes for interviewers to substitute personal preference.

\subsection*{10. Use a Common 5-Level Proficiency Scale}

The following is a \textbf{founder-designed rubric}, not a psychometrically validated scale. Its value is consistency; the exact labels are less important than concrete anchors.

\begin{table}[htbp]
\centering
\small
\caption{Common 5-Level Proficiency Scale}
\label{tab:proficiency_scale}
\begin{tabularx}{\columnwidth}{c p{2.2cm} X}
\toprule
\textbf{Score} & \textbf{Label} & \textbf{Observable Meaning} \\
\midrule
1 & Unsafe / not demonstrated & Major errors or no relevant evidence; requires continuous correction. \\
2 & Assisted & Can perform fragments, but misses important issues or needs repeated prompts. \\
3 & Independent in routine cases & Delivers acceptable ordinary work, flags exceptions, follows standards. \\
4 & Independent in difficult cases & Handles ambiguity and difficult cases, explains trade-offs, prevents recurring errors. \\
5 & Institution-building & Creates methods, teaches others, and improves standards beyond their own output. \\
\bottomrule
\end{tabularx}
\end{table}

For each role, \textbf{rewrite these anchors in occupational language}. A ``4'' in classroom teaching should describe observable diagnosis and adaptation; a ``4'' in backend engineering should describe system behavior under load or failure. Do not pretend the numerals are precise measurements of a person's intrinsic worth. Record the underlying evidence and uncertainty. A single exercise cannot reliably distinguish a 4.1 from a 4.3.

\subsection*{11. Classify Competencies by Readiness, Trainability, and Consequence}

A useful 4-cell map separates the competence that must exist at entry from what may develop on the job:
\begin{enumerate}[leftmargin=*]
\item \textbf{Entry-critical and hard to train quickly}: ethical judgment, baseline mathematical reasoning in a quantitative role, production safety, accurate reading of a learner's explanation, or legal subject competence where required.
\item \textbf{Entry-critical but trainable quickly}: an internal tool or documentation protocol when initial supervision is available.
\item \textbf{Future-critical and trainable}: a new cloud vendor, curriculum topic, or account vertical.
\item \textbf{Peripheral}: skills whose absence has little effect on the role's economic or quality objective.
\end{enumerate}

The classification depends on context. A new programming framework may be peripheral for an outstanding systems engineer but entry-critical for a 6-week client deployment with no mentoring. A theoretical subject might be learnable by a teacher over months but not before a class beginning tomorrow. Record the reason for each classification rather than treating any category as permanent.

\subsection*{12. Separate 4 Forms of Evidence}

\textbf{Past performance evidence} includes completed work, outcomes, decisions, failures, and references. \textbf{Direct performance evidence} comes from job-relevant work samples and simulations. \textbf{Learning evidence} comes from observing how a candidate incorporates correction and adapts to novel tasks. \textbf{Behavioral and integrity evidence} comes from consistent responses to real dilemmas, transparent limitations, and corroborated interactions. These sources have different weaknesses.

Portfolios can be selectively curated; take-homes can be externally assisted; interviews invite polished storytelling; short tests produce sampling error; references may be chosen for friendliness. More methods are not always better: 5 interviews all probing ``confidence'' are not 5 independent observations. The objective is \textbf{triangulation}, including evidence that could disconfirm the appealing story.

\subsection*{13. Build the Exercise Before the Finalist Is Known}

Assessment design should not be tailored to reward someone the founder already favors. Create the task, expected good and bad answers, scoring criteria, time box, permissible tools, and escalation rules 1st. Calibrate the rubric on an internal example, including a deliberately flawed answer. Where possible, anonymize outputs during initial scoring and have 2 independent evaluators score the most consequential tasks before discussing impressions.

A good exercise is \textit{small but complete}: a sales discovery call with an actual next-step decision; a technical debugging incident ending in a tested patch and explanation; a teaching segment followed by a misconception diagnosis; a research critique ending with a decision about whether the claim is supportable. Ask candidates to explain what they would \textbf{not} do as well as what they would do. The omissions reveal judgment.

\subsection*{14. A 2-Stage Assessment Is Usually Better Than 1 Heroic Test}

Stage 1 should be low-burden: evidence of core eligibility, a short structured screen, and a small role-relevant probe. Stage 2 may involve a paid work sample, teaching demonstration, paired debugging exercise, client simulation, or leadership case, followed by focused structured interviews and references. Reserve the highest-burden activities for finalists. Make the process, expected time, compensation, tool permissions, and data handling explicit.

For an early founder, an efficient pattern is:
\begin{center}
\textbf{30-minute structured screen $\to$ 60-120-minute paid sample $\to$ 40-minute evidence interview $\to$ 2 focused references $\to$ decision memo}.
\end{center}
An executive assessment or safety-critical position may justify much more; a casual part-time role may justify less. The principle is proportionality: assessment burden should track the expected value of better information and the costs imposed on candidates.
